\honest{Evolving to Extinction}{Eliezer Yudkowsky}{2007}

It is a common misconception, I say, that evolution works for the good of the species. Evolution acts on a population over time, but what changes are gene frequencies, and a gene spreads according to what it does for its own copies. \nb{This is the central claim of George C.~Williams's \textsc{Adaptation and Natural Selection} (1966), which the post does not name.}

My first example is the sex ratio. A group whose mothers had three daughters for every son would have 50\% more grandchildren than a group with equal numbers. But every child has one father and one mother, so the rarer males are, the more each son is worth to the parent. Individual selection therefore keeps investment in sons and daughters equal, and a maternity ward shows that it beats group selection in humans. \nb{The argument is usually called Fisher's principle. In species whose young mate among themselves, female-biased ratios are common; some biologists explain them by competition among brothers, others by group selection.}

A gene that sacrifices its carrier to save the whole species, which I call a Frodo gene, loses frequency, and if extinction threats recur, the species dies out. So it is ``perfectly possible'' for a species to evolve to extinction. My cases: viruses that kill their hosts too fast, which I guess ``probably happened any number of times''; a segregation distorter on the male sex chromosome of some mice, which makes only sons; transposons, which make up around half the maize genome and ``may not extinguish a species''; and a perfect DNA-copying mechanism, which would spread and then leave the species unable to adapt. \nb{Hamilton showed in 1967 that a driving Y chromosome can take a population extinct. But the mouse distorter that group selection was invoked for is the t haplotype, on chromosome 17, not a sex chromosome, and the editor could not find the mouse gene the post describes. A laboratory case of evolving to extinction did exist: cheating bacteria that drove their populations extinct (Fiegna and Velicer, 2003).}

Then the bystander effect. A student having a seizure was helped 85\% of the time by one bystander and 31\% of the time by five. ``I speculate'' that an arms race arose to be the last to step forward. Humanity faces extinction threats now, and ``there ain't a lot of people steppin' forward.'' \nb{The post gives no evidence for this last claim.}

Cancer cells outcompete their neighbours until the body dies, and bodies exist only by suppressing evolution among their cells. So do not praise evolution for its concern for individuals, species or even genes. ``Humanity may be finishing up the process right now.'' \nb{Of the cases of extinction in the post, only cancer is an observed one. The rest are thought experiments or hedged guesses.}
